\section{Introduction}


The recent advances in quantum information and quantum computation have profoundly transformed computational science.
There are many scientific fields that will potentially benefit from the existence of a full-scale quantum computer, 
including economics, physics (via efficient simulation of quantum systems)
and even machine learning (see e.g.~\cite{du2025quantummachinelearninghandson}).
However, the most significant and impactful application of quantum computing is in the field of
cybersecurity, as Shor's algorithm~\cite{Shor_1997} allows to break conventional asymmetric cryptography
based on prime factorization or discrete logs exponentially faster than any classical computer.
This discovery fundamentally disrupted cryptographic research, paving the way for the emergence of a new field
called post-quantum cryptography, in which new cryptographic systems are designed which conjecturally resist the
dangers of quantum algorithms.
Moreover, quantum phenomena such as superdense coding~\cite{Harrow_2004} and quantum teleportation~\cite{PhysRevLett.70.1895}
lays the groundwork for the field of quantum cryptography,
which is the study of quantum protocols resistant to eavesdropping.

A formal method to quantify these seemingly information-theoretic advantages of quantum physics
comes from the study of \emph{non-local games}.
A non-local game is an interactive cooperative game, in which the players win by fulfillling 
a given relation between the questions that are given to each player and their respective answers. 
While there is no communication allowed between the players during the game, possibly entangled quantum resources are allowed to be shared.
By leveraging this quantum advantage it is possible for the players to achieve winning
probabilities that exceed those of any classical (i.e. non-quantum) strategy.

Perhaps the best known non-local game is the magic square game
introduced in~\cite{aravind2003simpledemonstrationbellstheorem},
in which a $3 \times 3$ grid of $\pm1$ entries must be (partially) filled by the players,
such that each row multiplies to $1$ and each column multiplies to $-1$.
The first player gets asked the assignments of a specific row, while the second player
is asked about a column. The game is won if the mutual entry of the magic square is
filled with the same values by both players.
While any strategy that involves entirely filling the magic square beforehand and answering 
according to this assignment is impossible by design, it is possible to leverage the quantum resources to 
win the game with certainty.
Due to this seemingly magical phenomenon,
the magic square game is an example of what is called a \emph{pseudo-telepathic} game.

Another example of such a \emph{pseudo-telepathic} game is the graph isomorphism game,
which is parametrized by two publicly known graphs $G$ and $H$.
The referee’s questions to each player comprises of a vertex of $G$,
to which the players must answer with a vertex of $H$. 
In order to win, the nodes given as answers must fulfill the same adjacency relations in
$H$ as the nodes that are given as questions by the referee do in $G$.
While one can construct a perfect classical game strategy if and only if $G$ and $H$ are isomorphic, 
it is known that non-isomorphic graph pairs exist for which a perfect quantum strategy can be found.
Such graph pairs are said to be quantum-isomorphic.
